\setcounter{equation}{0}
\section{ Models of Computation}

Classical computers can be modeled by a Turing machine. A \emph{Turing Machine} is an abstract model of computation comprised of an infinitely long sequential tape, a read/write head, a set of symbols, and a set of rules for reading / writing and moving the head based on the symbol under it. These rules are called \emph{functions}, and the symbol under the head is called the \emph{state}. The machine moves from one state to another by executing a matching transition function. If no transition function matches, the machine halts.

The \emph{Church--Turing Thesis} \cite{ct36, tt36} says the following: a Turing machine can compute any function computable by a reasonable physical device. Programming languages can be translated into Turing machines and vice versa. Thus, the computation of any reasonable function can be expressed by a programming language.

The \emph{extended} (or \emph{strong}) Church--Turing thesis goes further than this and asserts that a probabilistic Turing machine can simulate any physical device with at most polynomial overhead. This is not a statement made by Church or Turing in their 1936 papers; it is a later thesis about the resources required for simulation. A \emph{probabilistic Turing Machine} can select among a set of transition functions with varying probability, instead of just one. In the following example, real numbers preceding a state indicate the probability of transitioning to that state.

\[ \begin{array}{ccccccc}
\ket{0} & \rightarrow & 0.5\ket{1} & + & 0.5\ket{2} \\
\ket{1} & \rightarrow & 0.2\ket{0} & + & 0.8\ket{1} \\
\ket{2} & \rightarrow & 0.4\ket{0} & + & 0.3\ket{2} & + & 0.3\ket{3} \\
\ket{3} & \rightarrow & 0.3\ket{0} & + & 0.7\ket{2} \end{array}\] 

Such machines can produce solutions with a negligible probability of error, which would be non computable if no error was tolerated. Probabilistic computation with negligible error is important for quantum computation.

